% !PRINTED 406
\weberpage{406}
Setzt man entsprechende Coëfficienten einander gleich, so
folgt
\[
\begin{aligned}
1.&\quad \alpha^2 - 2\beta = 36,\\
2.&\quad \beta^2 + 2\delta - 2\alpha\gamma = 528,\\
3.&\quad \gamma^2 + 2\alpha\varepsilon - 2\beta\delta = 3897,\\
4.&\quad \delta^2 - 2\gamma\varepsilon = 14354,\\
5.&\quad \varepsilon^2 = 21025,
\end{aligned}
\]
und es sind nun ganzzahlige Werthe von $\alpha, \beta, \gamma, \delta, \varepsilon$ zu suchen,
die diesen fünf Gleichungen genügen. Zunächst erhält man aus 5.:
\[
\varepsilon = \sqrt{21025} = 145 = 5\,.\,29;
\]
und man kann $\varepsilon$ positiv annehmen, weil das von $x$ unabhängige
Glied jedenfalls in einem der beiden Factoren positiv ist. Die
Gleichung 4. ergiebt, wenn man rechts den Rest nach dem Modul
145 nimmt,
\weq{4}{\delta^2 \equiv -1 \pmod{145},}
und diese Congruenz hat die vier Lösungen
\[
\delta \equiv \pm 12,\ \pm 17 \pmod{145}.
\]
Da $\delta$ nach 4. ausserdem gerade sein muss, so könnte $\delta$ folgende
Werthe haben:
\[
\delta = \pm 12,\ \pm 128,\ \pm 162.
\]
Weiter wird man vorläufig nicht gehen, da grosse Zahlwerthe für
$\delta$ von vornherein unwahrscheinlich sind. Diesen Werthen von $\delta$
entsprechend erhält man aus 4. für $\gamma$:
\[
\gamma = -49,\ 7,\ 41.
\]
Wenden wir noch den Modul 5 an, so folgt:
\[
\gamma^2 \equiv 1,\ -1,\ 1 \pmod{5}
\]
und aus 3.:
\[
\beta \equiv \pm 1,\ \pm 2,\ \pm 1 \pmod{5}
\]
und aus 1.:
\[
\alpha^2 \equiv 1 \pm 2,\ 1 \mp 1,\ 1 \pm 2 \pmod{5}.
\]
Daraus ergiebt sich, da nur 0, 1 und 4, nicht 2 und 3 nach
dem Modul 5 mit einem Quadrat congruent sein können, dass
in $\delta$ die Zeichen so genommen werden müssen:
\[
\begin{aligned}
\delta &= -12,\ +128,\ -162\\
\alpha &\equiv \pm 2,\ 0,\ \pm 2 \pmod{5}.
\end{aligned}
\]
Die Gleichung 2. zeigt noch, dass der mittlere Fall auszuschliessen ist, da die rechte Seite nicht durch 5 theilbar ist.
